\par
\phantomsection\addcontentsline{toc}{chapter}{原课程大纲与考核}
\noindent{\color{structurecolor}\bfseries\fontsize{14.4}{20}\selectfont 原课程大纲与考核}\par
\textbf{课程.} MIT 18.335J Introduction to Numerical Methods (数值方法导论), 2019 年春季, 研究生课程. 任课教师为 Steven G. Johnson; 第 7、8 讲的 SVD/GSVD 为 Alan Edelman 客座讲授. 课程属数学系及电气工程与计算机科学系. 官方公布的授课频率为每周 3 次、每次 1 小时; OCW 大纲只给频率; 同学期 MIT 教师课程仓库的 \texttt{spring19} 分支另明确注明周一、周三、周五 15:00--16:00, 教室 2-190; 答疑为周四 16:00--17:00, 2-345. 时间按 MIT 当地钟点理解, 不据历史试卷推定.

\textbf{先修与工具.} 18.06、18.700 线性代数或同等基础. 作业使用 Julia, 无需已有 Julia 经验. 指定教材为 Trefethen 与 Bau 的 \emph{Numerical Linear Algebra} (SIAM, 1997). 本册正文按数值线性代数重组, 作业和考试保持原编号及年份.

\textbf{原大纲.} 线性系统的直接法与迭代法、特征值分解、QR/SVD、数值稳定性与精度、IEEE 浮点标准、稀疏及结构矩阵和线性代数软件; 还可能包括存储层次与缓存、非线性优化、数值积分、FFT 和灵敏度分析. 后几个专题在本系列其他科目中展开.

\textbf{考核.} 四份作业占 33\%, 带回家完成的期中考试占 33\%, 期末项目占 34\%; 没有另列期末笔试. 项目为 8--15 页、单栏单倍行距的综述论文, 介绍课上未讲的数值算法, 给出自己运行的数值结果, 分析理论及实验精度和性能, 公平比较至少一个竞争算法, 并在正文充分引用文献. 可以把 PDE 离散矩阵作为输入研究求解器, 但项目主题不应是 PDE 离散技术本身. 提案为约一页的非正式计划及若干起始文献.

\textbf{课程协作规则.} 可讨论和阅读资料, 但不得参考往年作业答案; 应先独立认真尝试, 再讨论或检索; 最后必须不看他人答案独立写出解答. 当期期中允许自己的课程笔记、教材及已发布课程资料, 禁止其他材料和互联网, 范围为第 20 讲及以前和 PS4. 本册为已公开课程材料的课后学习整理, 不改变这些原规则.

以下是 \href{https://ocw.mit.edu/courses/18-335j-introduction-to-numerical-methods-spring-2019/pages/calendar/}{2019 年官方课历}的中文对应. 原表按周和讲次给出关键节点, 没有绝对日期; 历史试卷的时限另在各卷保留, 不合并进本表. 同年课程仓库的绝对日期另列在本表后, 两个来源的讲次与截止表述存在差别, 不强行合并.
\begin{longtable}{p{.07\linewidth}p{.62\linewidth}p{.22\linewidth}}
\toprule 周 & 讲次与主题 & 考核节点\\\midrule\endhead
1 & 1 课程概览、Newton 求根; 2 浮点运算 & \\
2 & 3 浮点求和与后向稳定; 4 向量空间范数; 5 条件数 & PS1 截止\\
3 & 6 常微分方程; 7 SVD、应用与条件数; 8 线性回归与 GSVD & \\
4 & 9 正规方程、QR 与 Gram--Schmidt; 10 MGS、Householder QR; 11 矩阵运算、缓存与分块 & PS2 截止\\
5 & 12 缓存无关算法、空间局部性; 13 LU、部分主元; 14 Cholesky、特殊求解器、Schur & \\
6 & 15 伴随矩阵、病态性、Hessenberg; 16 幂法、QR; 17 移位 QR、Rayleigh 商 & PS3 截止\\
7 & 18 Krylov、Arnoldi; 19 Arnoldi/Lanczos 重启; 20 GMRES 及收敛 & 项目提案截止\\
8 & 21 预条件与 CG; 22 CG 收敛; 23 BiCG & PS4 截止\\
9 & 24 稀疏直接法; 25 优化概览; 26 伴随方法 & 带回家期中在第 25 讲前截止\\
10 & 27 特征值与递推的伴随方法; 28 信赖域与 CCSA & \\
11 & 29 Lagrange 对偶; 30 拟 Newton 与 BFGS; 31 BFGS 推导 & \\
12 & 32 无导数优化; 33 积分与梯形法; 34 Clenshaw--Curtis & \\
13 & 35 Chebyshev 逼近; 36 带权积分与 Gauss 求积; 37 自适应、多维求积 & \\
14 & 38 DFT 与 FFT; 39 FFT 与 FFTW & 期末项目在学期末截止\\\bottomrule
\end{longtable}
\textbf{同年教师课程仓库的日期记录.} \href{https://github.com/mitmath/18335/blob/spring19/README.md}{MIT \texttt{spring19} 分支}另列 PS1 为 2 月 15 日、PS2 为 3 月 1 日 15:00、PS3 为 3 月 15 日、PS4 为 4 月 5 日; 项目提案为 3 月 22 日, 期中于 4 月 8 日发布、4 月 9 日截止, 项目于 5 月 15 日截止. 该 README 的计划写作业约六份, 而 OCW 大纲与资源索引实际公布四份. 它把第 25 讲记为 4 月 8 日, 与 OCW 表的“期中在第 25 讲前截止”不能直接拼成同一绝对时间表; 本册保持各来源的原陈述, 四份作业统计以实际 OCW 题纸为准. 历史 2008--2013、2015 试卷不共享这些日期.

\textbf{考核材料与缺口.} 四份作业共 15 道大题. 当期及 2008--2013、2015 年历史期中共 29 道数学大题、60 个子问; 2019 年 Problem 0 是诚信声明. 作业纸共列 29 个题目/子问位置, 其中 11 处只引用外部教材习题, 原纸没有完整题面. 本册保留这些题号, 翻译并审校 MIT 官方解答中实际给出的结论, 不将它们冒称教材原题全文. 教材内未见的子问不计入统计. 2009 年无官方解答, 本册相应解答明确标为 AI 补写. PS1 第 2、3 题的完整官方解答来自官方链接的 notebook, 答案 PDF 只作指引. 具体来源、逐题冻结与文件哈希见随源码的清单.

来源: \href{https://ocw.mit.edu/courses/18-335j-introduction-to-numerical-methods-spring-2019/pages/syllabus/}{课程大纲}, \href{https://ocw.mit.edu/courses/18-335j-introduction-to-numerical-methods-spring-2019/pages/resource-index/}{资源索引}, \href{https://ocw.mit.edu/courses/18-335j-introduction-to-numerical-methods-spring-2019/pages/week-9/}{期中规则与历年试卷}. 作者及许可沿用本册前言和来源附录.
